% !TEX root = ../thesis.tex

\chapter{AI Tools}\label{app:ai_tools}

This appendix provides a transparent overview of the artificial intelligence and \gls{ai}-assisted tools that were used during the preparation of this thesis. The purpose of this appendix is not to present these tools as part of the proposed methodological contribution, but to clearly distinguish between the author's own work and external \gls{ai}-assisted support used during literature search, workflow design, implementation support, debugging, and text refinement.

The use of \gls{ai}-assisted tools was limited to supportive, consultative, and organizational tasks. These tools were not used as autonomous decision-making systems and did not replace the author's responsibility for the research design, implementation, experimental evaluation, interpretation of results, or final academic writing. All methodological choices, experimental decisions, code integration, result interpretation, and final text approval remained under the author's control.

\section{General Principles of Use}

\gls{ai}-assisted tools were used according to the following principles:

\begin{itemize}
    \item \gls{ai} tools were used as auxiliary instruments for literature discovery, note enrichment, workflow consultation, code assistance, debugging, and language refinement.
    \item The tools did not independently define the research problem, objectives, experimental protocol, or final conclusions of the thesis.
    \item Generated code snippets were reviewed, adapted, and integrated manually by the author before being used in the project.
    \item Interpretations of experimental results were treated as suggestions and were checked against the measured metrics, generated visualizations, and the broader methodological context of the thesis.
    \item \gls{ai}-assisted textual edits were used for improving clarity, grammar, structure, and consistency, while the final meaning, argumentation, and responsibility for the text remained with the author.
    \item The tools were not used to fabricate sources, results, measurements, tables, or figures.
\end{itemize}

\section{ResearchRabbit}

ResearchRabbit was used as a literature discovery tool. It supported the search for relevant scientific articles and made it possible to explore citation networks between already identified publications. This was useful for detecting closely related works, following citation paths, and identifying papers that were methodologically or conceptually connected to the thesis topic.

In addition, ResearchRabbit helped trace the origin and usage context of several domain-specific terms encountered during the literature review. However, the tool was not used to determine the final relevance of a publication automatically. The final selection of sources was performed manually by reading, screening, and comparing the articles in the context of the thesis topic.

\section{NotebookLM}

NotebookLM was used for the preliminary analysis of a collected cluster of scientific articles. Its use started only after manual screening had already been performed. The tool helped summarize and compare selected papers and supported the identification of articles that were likely to be relevant for the analytical part of the thesis.

NotebookLM was also used to supplement manually created notes for individual articles. This was done in order to reduce the risk of overlooking details that could be important in the context of \gls{semanticsegmentation}, \gls{lidarderiveddata} visualizations, archaeological object detection, label uncertainty, \gls{classimbalance}, or evaluation methodology.

The tool was additionally used when preparing indirect citations. In such cases, NotebookLM helped identify whether a statement, although not explicitly formulated in a given article, was indirectly supported by the article's results, methodology, or discussion. These suggestions were not accepted automatically. The author verified the context manually before using such material in the thesis.

\section{Anara}

Anara was used only for the initial analysis of the root article by Bundzel et al., which served as one of the main starting points for understanding the domain-specific background of \gls{lidar}-based archaeological segmentation. Its use was limited to early familiarization with the article and did not replace the author's own reading, interpretation, or subsequent analysis of the publication.

\section{ChatGPT 5.4 Thinking Extended}

ChatGPT 5.4 Thinking Extended was used mainly as a consultative tool during the design of a reproducible research workflow. It supported the planning of a project structure that separated source code, data, configurations, experiment outputs, evaluation results, and thesis-related artifacts. This helped reduce the amount of time spent on organizational and implementation-related decisions and allowed greater focus on experimental design and result interpretation.

The tool was also useful during the initial analysis of experimental results. It helped formulate possible interpretations of measured metrics, compare experiment groups, and identify relationships between model architecture, \gls{inputrepresentation}, label representation, sampling strategy, \gls{lossfunction}, and observed performance.

In addition, ChatGPT 5.4 Thinking Extended was used to generate small code snippets for specific functions, pandas data frames for comparative result tables, and helper scripts for processing intermediate experiment outputs. It was also used for targeted debugging of code. All generated code was reviewed, modified, and tested by the author before being included in the workflow.

\section{ChatGPT 5.5 Thinking Extended}

ChatGPT 5.5 Thinking Extended was used during the later stages of thesis preparation. It supported discussion of experimental results, checking consistency between different sections of the thesis, identifying possible grammar issues, and improving the formulation of selected paragraphs.

The tool was also used to assist with the presentation of results, including the formatting of tables, graphs, figures, and explanatory text. Its role was editorial and consultative. It did not replace the author's responsibility for the results, interpretation, structure, or final wording of the thesis.

\section{Boundaries of AI Assistance}

The use of \gls{ai}-assisted tools did not include the following:

\begin{itemize}
    \item autonomous formulation of the thesis objectives;
    \item autonomous selection of the final methodology;
    \item independent implementation of the complete experimental pipeline;
    \item automatic generation of experimental results;
    \item fabrication or invention of citations, metrics, tables, or figures;
    \item replacement of manual validation of code, results, and written claims;
    \item replacement of the author's interpretation of the archaeological and machine-learning relevance of the results.
\end{itemize}

The final thesis therefore remains the author's own academic work. \gls{ai}-assisted tools were used only as supporting instruments for literature exploration, workflow organization, implementation assistance, debugging, language refinement, and critical discussion of results.